% !PRINTED 458
\weberpage{458}
Der Satz ist zunächst evident, wenn die Functionen $\Phi_1, \Phi_2, \Phi_3\ldots$ nur von einer Veränderlichen abhängen; denn dann giebt es überhaupt nur eine endliche Anzahl von Zahlwerthen für diese Veränderliche, wofür eine dieser Functionen Null wird.

Dann aber können wir die Richtigkeit des Satzes für Functionen von $n+1$ Veränderlichen leicht einsehen, falls wir ihn für Functionen von $n$ Veränderlichen als bewiesen betrachten. Denn ordnen wir die Functionen nach der $(n+1)^{\text{ten}}$ Veränderlichen $t$, so können wir für die übrigen $n$ Veränderlichen nach Voraussetzung solche rationale Werthe setzen, dass in keiner der Functionen alle Coëfficienten der Potenzen von $t$ verschwinden; dann haben wir Functionen der einen Veränderlichen $t$ und können für diese einen solchen rationalen Werth setzen, dass keine der Functionen verschwindet. Damit ist also der Satz allgemein bewiesen.

Es soll jetzt zunächst nachgewiesen werden, dass man die gleichzeitige Adjunction \gesperrt{mehrerer} algebraischen Grössen durch die Adjunction \gesperrt{einer einzigen} ersetzen kann, mit anderen Worten, dass jeder Körper $\Omega(\alpha, \beta, \gamma\ldots)$ angesehen werden kann als ein Körper $\Omega(\alpha)$.

Es seien also $\alpha, \beta, \gamma\ldots$ algebraische Grössen in beliebiger Anzahl und

\weq{1}{A(x) = 0, \quad B(x) = 0, \quad C(x) = 0\ldots}

Gleichungen in $\Omega$, deren Wurzeln $x=\alpha$, $x=\beta$, $x=\gamma\ldots$ sind. Keine dieser Gleichungen soll eine mehrfache Wurzel haben, eine Voraussetzung, durch die die Allgemeinheit nicht beschränkt wird.

Wir haben im vorigen Paragraphen gesehen, dass sich jede Grösse eines Körpers $\Omega(\alpha)$ als \gesperrt{ganze rationale Function} in $\Omega$ von $\alpha$ darstellen lässt. Setzen wir in diesem Satze an Stelle des Körpers $\Omega$ den Körper $\Omega(\beta, \gamma\ldots)$, so folgt, dass jede Zahl in $\Omega(\alpha, \beta, \gamma\ldots)$ als ganze rationale Function von $\alpha$ mit Coëfficienten aus $\Omega(\beta, \gamma\ldots)$ dargestellt werden kann; und wenn wir dieselbe Schlussweise in Beziehung auf die Coëfficienten wiederholen, so ergiebt sich:

2. \gesperrt{Jede Grösse des Körpers $\Omega(\alpha, \beta, \gamma\ldots)$ kann als ganze rationale Function in $\Omega$ von $\alpha, \beta, \gamma\ldots$ dargestellt werden.}
